\section{Associated points}
\label{section-associated}

\noindent
Let $R$ be a ring and let $M$ be an $R$-module.
Recall that a prime $\mathfrak p \subset R$ is {\it associated} to $M$
if there exists an element of $M$ whose annihilator is $\mathfrak p$.
See Algebra, Definition \ref{algebra-definition-associated}.
Here is the definition of associated points
for quasi-coherent sheaves on schemes
as given in \cite[IV Definition 3.1.1]{EGA}.

\begin{definition}
\label{definition-associated}
Let $X$ be a scheme.
Let $\mathcal{F}$ be a quasi-coherent sheaf on $X$.
\begin{enumerate}
\item We say $x \in X$ is {\it associated} to $\mathcal{F}$
if the maximal ideal
$\mathfrak m_x$ is associated to the $\mathcal{O}_{X, x}$-module
$\mathcal{F}_x$.
\item We denote $\text{Ass}(\mathcal{F})$ or $\text{Ass}_X(\mathcal{F})$
the set of associated points of $\mathcal{F}$.
\item The {\it associated points of $X$} are the associated
points of $\mathcal{O}_X$.
\end{enumerate}
\end{definition}

\noindent
These definitions are most useful when $X$ is locally Noetherian
and $\mathcal{F}$ of finite type.
For example it may happen that a generic point of an irreducible
component of $X$ is not associated to $X$, see
Example \ref{example-no-associated-prime}.
In the non-Noetherian case it may be more convenient to use weakly
associated points, see
Section \ref{section-weakly-associated}.
Let us link the scheme theoretic notion with the algebraic notion
on affine opens; note that this correspondence works perfectly only
for locally Noetherian schemes.

\begin{lemma}
\label{lemma-associated-affine-open}
Let $X$ be a scheme. Let $\mathcal{F}$ be a quasi-coherent sheaf on $X$.
Let $\Spec(A) = U \subset X$ be an affine open, and set
$M = \Gamma(U, \mathcal{F})$.
Let $x \in U$, and let $\mathfrak p \subset A$ be the corresponding prime.
\begin{enumerate}
\item If $\mathfrak p$ is associated to $M$, then $x$ is associated
to $\mathcal{F}$.
\item If $\mathfrak p$ is finitely generated, then the converse holds
as well.
\end{enumerate}
In particular, if $X$ is locally Noetherian, then the equivalence
$$
\mathfrak p \in \text{Ass}(M) \Leftrightarrow x \in \text{Ass}(\mathcal{F})
$$
holds for all pairs $(\mathfrak p, x)$ as above.
\end{lemma}

\begin{proof}
This follows from
Algebra, Lemma \ref{algebra-lemma-associated-primes-localize}.
But we can also argue directly as follows.
Suppose $\mathfrak p$ is associated to $M$.
Then there exists an $m \in M$ whose annihilator is $\mathfrak p$.
Since localization is exact we see that
$\mathfrak pA_{\mathfrak p}$ is the annihilator of
$m/1 \in M_{\mathfrak p}$. Since $M_{\mathfrak p} = \mathcal{F}_x$
(Schemes, Lemma \ref{schemes-lemma-spec-sheaves})
we conclude that $x$ is associated to $\mathcal{F}$.

\medskip\noindent
Conversely, assume that $x$ is associated to $\mathcal{F}$,
and $\mathfrak p$ is finitely generated.
As $x$ is associated to $\mathcal{F}$
there exists an element $m' \in M_{\mathfrak p}$ whose
annihilator is $\mathfrak pA_{\mathfrak p}$. Write
$m' = m/f$ for some $f \in A$, $f \not \in \mathfrak p$.
The annihilator $I$ of $m$ is an ideal of $A$ such that
$IA_{\mathfrak p} = \mathfrak pA_{\mathfrak p}$. Hence
$I \subset \mathfrak p$, and $(\mathfrak p/I)_{\mathfrak p} = 0$.
Since $\mathfrak p$ is finitely generated,
there exists a $g \in A$, $g \not \in \mathfrak p$ such that
$g(\mathfrak p/I) = 0$. Hence the annihilator of $gm$ is
$\mathfrak p$ and we win.

\medskip\noindent
If $X$ is locally Noetherian, then $A$ is Noetherian
(Properties, Lemma \ref{properties-lemma-locally-Noetherian})
and $\mathfrak p$ is always finitely generated.
\end{proof}

\begin{lemma}
\label{lemma-ass-support}
Let $X$ be a scheme.
Let $\mathcal{F}$ be a quasi-coherent $\mathcal{O}_X$-module.
Then $\text{Ass}(\mathcal{F}) \subset \text{Supp}(\mathcal{F})$.
\end{lemma}

\begin{proof}
This is immediate from the definitions.
\end{proof}

\begin{lemma}
\label{lemma-ses-ass}
Let $X$ be a scheme.
Let $0 \to \mathcal{F}_1 \to \mathcal{F}_2 \to \mathcal{F}_3 \to 0$
be a short exact sequence of quasi-coherent sheaves on $X$.
Then
$\text{Ass}(\mathcal{F}_2) \subset
\text{Ass}(\mathcal{F}_1) \cup \text{Ass}(\mathcal{F}_3)$
and
$\text{Ass}(\mathcal{F}_1) \subset \text{Ass}(\mathcal{F}_2)$.
\end{lemma}

\begin{proof}
For every point $x \in X$ the sequence of stalks
$0 \to \mathcal{F}_{1, x} \to \mathcal{F}_{2, x} \to \mathcal{F}_{3, x} \to 0$
is a short exact sequence of $\mathcal{O}_{X, x}$-modules.
Hence the lemma follows from
Algebra, Lemma \ref{algebra-lemma-ass}.
\end{proof}

\begin{lemma}
\label{lemma-finite-ass}
Let $X$ be a locally Noetherian scheme.
Let $\mathcal{F}$ be a coherent $\mathcal{O}_X$-module.
Then $\text{Ass}(\mathcal{F}) \cap U$ is finite for
every quasi-compact open $U \subset X$.
\end{lemma}

\begin{proof}
This is true because the set of associated primes of a finite module over
a Noetherian ring is finite, see
Algebra, Lemma \ref{algebra-lemma-finite-ass}.
To translate from schemes to algebra use that $U$ is a finite union of
affine opens, each of these opens is the spectrum of a Noetherian ring
(Properties, Lemma \ref{properties-lemma-locally-Noetherian}),
$\mathcal{F}$ corresponds to a finite module over this ring
(Cohomology of Schemes, Lemma \ref{coherent-lemma-coherent-Noetherian}),
and finally use
Lemma \ref{lemma-associated-affine-open}.
\end{proof}

\begin{lemma}
\label{lemma-ass-zero}
Let $X$ be a locally Noetherian scheme. Let $\mathcal{F}$ be a
quasi-coherent $\mathcal{O}_X$-module. Then
$$
\mathcal{F} = 0 \Leftrightarrow \text{Ass}(\mathcal{F}) = \emptyset.
$$
\end{lemma}

\begin{proof}
If $\mathcal{F} = 0$, then $\text{Ass}(\mathcal{F}) = \emptyset$
by definition. Conversely, if $\text{Ass}(\mathcal{F}) = \emptyset$,
then $\mathcal{F} = 0$ by
Algebra, Lemma \ref{algebra-lemma-ass-zero}.
To translate from schemes to algebra, restrict to any affine and use
Lemma \ref{lemma-associated-affine-open}.
\end{proof}

\begin{example}
\label{example-no-associated-prime}
Let $k$ be a field. The ring $R = k[x_1, x_2, x_3, \ldots]/(x_i^2)$
is local with locally nilpotent maximal ideal $\mathfrak m$.
There exists no element of $R$ which has annihilator $\mathfrak m$.
Hence $\text{Ass}(R) = \emptyset$, and $X = \Spec(R)$
is an example of a scheme which has no associated points.
\end{example}

\begin{lemma}
\label{lemma-restriction-injective-open-contains-ass}
Let $X$ be a locally Noetherian scheme. Let $\mathcal{F}$ be a quasi-coherent
$\mathcal{O}_X$-module. If $U \subset X$ is open and
$\text{Ass}(\mathcal{F}) \subset U$, then
$\Gamma(X, \mathcal{F}) \to \Gamma(U, \mathcal{F})$ is injective.
\end{lemma}

\begin{proof}
Let $s \in \Gamma(X, \mathcal{F})$ be a section which restricts to zero on $U$.
Let $\mathcal{F}' \subset \mathcal{F}$ be the image of the map
$\mathcal{O}_X \to \mathcal{F}$ defined by $s$. Then
$\text{Supp}(\mathcal{F}') \cap U = \emptyset$. On the other hand,
$\text{Ass}(\mathcal{F}') \subset \text{Ass}(\mathcal{F})$
by Lemma \ref{lemma-ses-ass}. Since also
$\text{Ass}(\mathcal{F}') \subset \text{Supp}(\mathcal{F}')$
(Lemma \ref{lemma-ass-support}) we conclude
$\text{Ass}(\mathcal{F}') = \emptyset$.
Hence $\mathcal{F}' = 0$ by Lemma \ref{lemma-ass-zero}.
\end{proof}

\begin{lemma}
\label{lemma-minimal-support-in-ass}
Let $X$ be a locally Noetherian scheme.
Let $\mathcal{F}$ be a quasi-coherent $\mathcal{O}_X$-module.
Let $x \in \text{Supp}(\mathcal{F})$ be a point in the support
of $\mathcal{F}$ which is not a specialization of another point of
$\text{Supp}(\mathcal{F})$. Then $x \in \text{Ass}(\mathcal{F})$.
In particular, any generic point of an irreducible component of $X$
is an associated point of $X$.
\end{lemma}

\begin{proof}
Since $x \in \text{Supp}(\mathcal{F})$ the module $\mathcal{F}_x$
is not zero. Hence
$\text{Ass}(\mathcal{F}_x) \subset \Spec(\mathcal{O}_{X, x})$
is nonempty by
Algebra, Lemma \ref{algebra-lemma-ass-zero}.
On the other hand, by assumption
$\text{Supp}(\mathcal{F}_x) = \{\mathfrak m_x\}$.
Since
$\text{Ass}(\mathcal{F}_x) \subset \text{Supp}(\mathcal{F}_x)$
(Algebra, Lemma \ref{algebra-lemma-ass-support})
we see that $\mathfrak m_x$ is associated to $\mathcal{F}_x$
and we win.
\end{proof}

\noindent
The following lemma is the analogue of
More on Algebra, Lemma \ref{more-algebra-lemma-check-injective-on-ass}.

\begin{lemma}
\label{lemma-check-injective-on-ass}
Let $X$ be a locally Noetherian scheme. Let
$\varphi : \mathcal{F} \to \mathcal{G}$ be a map of
quasi-coherent $\mathcal{O}_X$-modules.
Assume that for every $x \in X$
at least one of the following happens
\begin{enumerate}
\item $\mathcal{F}_x \to \mathcal{G}_x$ is injective, or
\item $x \not \in \text{Ass}(\mathcal{F})$.
\end{enumerate}
Then $\varphi$ is injective.
\end{lemma}

\begin{proof}
The assumptions imply that $\text{Ass}(\Ker(\varphi)) = \emptyset$
and hence $\Ker(\varphi) = 0$ by Lemma \ref{lemma-ass-zero}.
\end{proof}

\begin{lemma}
\label{lemma-check-isomorphism-via-depth-and-ass}
Let $X$ be a locally Noetherian scheme. Let
$\varphi : \mathcal{F} \to \mathcal{G}$ be a map of
quasi-coherent $\mathcal{O}_X$-modules. Assume $\mathcal{F}$ is coherent
and that for every $x \in X$ one of the following happens
\begin{enumerate}
\item $\mathcal{F}_x \to \mathcal{G}_x$ is an isomorphism, or
\item $\text{depth}(\mathcal{F}_x) \geq 2$ and
$x \not \in \text{Ass}(\mathcal{G})$.
\end{enumerate}
Then $\varphi$ is an isomorphism.
\end{lemma}

\begin{proof}
This is a translation of More on Algebra, Lemma
\ref{more-algebra-lemma-check-isomorphism-via-depth-and-ass}
into the language of schemes.
\end{proof}




